\begingroup
\thispagestyle{empty}
{\titlefont\LARGE چکیده\par}
\vspace{1.5em}
پردازش گفتار شاخه‌ای کاربردی از پردازش سیگنال و هوش مصنوعی است که به تحلیل، درک و تولید گفتار انسان می‌پردازد. این جزوه با معرفی وظایف پایین‌دستی در پردازش زبان، گفتار و بینایی و ترکیب آن‌ها در محصولات واقعی آغاز می‌شود. سپس تصویر کلی علوم زبان و آواشناسی تولیدی را مرور می‌کند و به آواشناسی آکوستیکی می‌رسد: تبدیل فوریه، هارمونیک‌ها، اسپکتروگرام، مدل منبع-فیلتر و فرمنت‌ها. در ادامه، کار عملی با نرم‌افزار \eng{Praat} و فضای واکه‌ای، ویژگی‌های الهام‌گرفته از شنوایی انسان مانند مقیاس مل و ضرایب \eng{MFCC}، و ویژگی‌های پروسودیک گفتار بررسی می‌شوند. فصل‌های پایانی به سامانه‌های تبدیل متن به گفتار می‌پردازند: نرمال‌سازی متن، رفع ابهام کلمات هم‌نگاشت، تبدیل نویسه به واج، مدل آکوستیکی و وکودر، و نیز بازنمایی متن با \eng{embedding}های ایستا و زمینه‌مند و افزودن اطلاعات ترتیب به آن‌ها با کدگذاری موقعیت.

\vfill
\textbf{کلیدواژه‌ها:} پردازش گفتار، آواشناسی، اسپکتروگرام، فرمنت، \eng{MFCC}، پروسودی، تبدیل متن به گفتار، \eng{embedding} زمینه‌مند، کدگذاری موقعیت
\clearpage
\endgroup

\begingroup
\thispagestyle{empty}
\begin{latin}
{\LARGE\bfseries Abstract\par}
\vspace{1.5em}
Speech processing is an applied branch of signal processing and artificial intelligence concerned with analysing, understanding and generating human speech. This booklet opens with downstream tasks in language, speech and vision, and shows how they combine into real products. It then reviews the big picture of the language sciences and articulatory phonetics, and moves on to acoustic phonetics: the Fourier transform, harmonics, spectrograms, the source--filter model and formants. Next come hands-on work with Praat and the vowel space, features inspired by human hearing such as the Mel scale and MFCCs, and the prosodic features of speech. The final chapters cover text-to-speech systems: text normalisation, homograph disambiguation, grapheme-to-phoneme conversion, acoustic models and vocoders, and text representation with static and contextual embeddings, together with positional encoding for injecting word order.

\vfill
\textbf{Keywords:} speech processing, phonetics, spectrogram, formants, MFCC, prosody, text-to-speech, contextual embeddings, positional encoding
\end{latin}
\clearpage
\endgroup

\begingroup
\thispagestyle{empty}
{\titlefont\LARGE دربارهٔ این پیش‌نویس\par}
\vspace{1em}
این نسخهٔ فشردهٔ جزوهٔ درس پردازش گفتار است و بر پایهٔ یادداشت‌های صوتی جلسات و تختهٔ کلاس تهیه شده است. در آن مطالبی که در کلاس بیشتر روی آن‌ها تمرکز شد و برای یادگیری و امتحان مهم‌ترند نگه داشته شده‌اند و توضیح‌های تکراری و فرعی حذف شده‌اند؛ برای جزئیات کامل‌تر به نسخهٔ مفصل جزوه مراجعه کنید. چون درس هنوز ادامه دارد، فقط جلسات اول تا دهم را پوشش می‌دهد. هر جا مطلبی خارج از گفته‌های کلاس اضافه شده باشد، در کادر «محتوای افزوده» آمده است.

\vspace{0.6em}
\textbf{راهنمای کادرها:}

\begin{definitionbox}{تعریف}
تعریف‌ها و مفاهیم کلیدی درس.
\end{definitionbox}
\begin{exambox}{نکتهٔ استاد / امتحان}
نکته‌هایی که استاد بر آن‌ها تأکید کرده است، اطلاعات ارزشیابی و امتحان، تمرین‌ها و نکته‌های جانبی کلاس.
\end{exambox}
\begin{correctionbox}{اصلاح یا دقت علمی}
جاهایی که گفته‌های کلاس یا یادداشت‌ها نیاز به اصلاح یا دقیق‌تر شدن داشته‌اند.
\end{correctionbox}
\begin{warningbox}{هشدار / خطای رایج}
اشتباه‌های رایج و نکته‌هایی که باید با احتیاط به کار بروند.
\end{warningbox}
\begin{assignmentbox}{تمرین و پروژه}
تمرین‌ها و فعالیت‌های عملی که در کلاس مطرح شده‌اند.
\end{assignmentbox}
\begin{addedbox}{مثال یا توضیح}
مثال‌های عددی، شکل‌ها و توضیح‌هایی که در کلاس گفته نشده‌اند و برای فهم بهتر اضافه شده‌اند.
\end{addedbox}
\clearpage
\endgroup
